\documentclass[11pt]{article}
\usepackage[T1]{fontenc}
\usepackage{lmodern,amsmath,amssymb,amsthm,geometry,hyperref}
\geometry{margin=1in}
\hypersetup{colorlinks=true,linkcolor=black,urlcolor=blue,citecolor=black}
\newtheorem{theorem}{Theorem}[section]
\newtheorem{lemma}[theorem]{Lemma}
\newtheorem{proposition}[theorem]{Proposition}
\newcommand{\Q}{\mathbf Q}\newcommand{\R}{\mathbf R}\newcommand{\Z}{\mathbf Z}
\newcommand{\C}{\mathbf C}\newcommand{\eps}{\varepsilon}\newcommand{\dd}{\mathrm d}
\newcommand{\e}{\mathrm e}\newcommand{\im}{\operatorname{Im}}\newcommand{\re}{\operatorname{Re}}
\newcommand{\dist}{\operatorname{dist}}\newcommand{\av}{\operatorname{Av}}
\newcommand{\ord}{\operatorname{ord}}\newcommand{\Gal}{\operatorname{Gal}}
\newcommand{\FF}{\mathbf F}\newcommand{\A}{\mathcal A}\newcommand{\Lfun}{\mathcal L}
\allowdisplaybreaks
\title{Cyclotomic logarithms, Kummer traces and modular-symbol quadrature for a second central derivative}
\author{\small\texttt{admin@cosmosgeek.org}\\[-1pt]\small\texttt{b321014@iiit-bh.ac.in}}
\date{}
\begin{document}
\maketitle
\begin{abstract}
For a weight-two elliptic newform whose central value and first derivative vanish,
we give a finite sum at squarefree modular-symbol conductor \(bq\). Its logarithms
are those of global cyclotomic units. We compute its first finite-dilation
correction, including both \(L(f,3)\) and \(L'(f,3)\). A second construction
uses Weil functions of two independent rational points. Along an explicit choice of
scales, a single complex embedding and a normalized average over Kummer
translations both recover the central second derivative. Under a two-prime
Kurihara unit hypothesis, a third test gives bounded integral weights
with prescribed residues and the same normalized real limit. A control
on the same elliptic curve retains a unit residue while tending to zero.
The proofs use a
modular-symbol distribution estimate, two beta integrals and the finite
multiplication formulas for the logarithm and cotangent.
\end{abstract}

\section{The squarefree modular-symbol formula}
Let \(E/\Q\) have conductor \(N\), with normalized newform
\(f(z)=\sum_{n\ge1}a_n\e^{2\pi inz}\). Write \(L(s)=L(E,s)\), in the
normalization with center \(1\). Set
\[
 F(z)=\sum_{n\ge1}\frac{a_n}{n}\e^{2\pi inz}
     =2\pi i\int_{i\infty}^{z}f(w)\,\dd w .
\]
At rational cusps use the limiting Eichler integral. Fix positive real period
scales \(\Omega_+,\Omega_-\), and put
\(s_+(x)=\re F(x)/\Omega_+\), \(s_-(x)=\im F(x)/\Omega_-\).
Thus a purely imaginary minus period has been replaced by its positive
magnitude. The choice of scales appears identically on the two sides of
every formula. Throughout assume \(L(1)=L'(1)=0\). Write
\(d(x)=2|\sin\pi x|\), \(\|x\|=\dist(x,\Z)\), and
\(\zeta_h=\exp(2\pi i/h)\).

Fix a prime \(b\nmid N\). Let \(n=b^j\), \(q\ne b\) a prime prime to \(N\),
\(L_0=bq,\ M=nL_0\) and \(H=\lceil q^{1/8}\rceil\).
For \(r=0,1\) define
\[
 S_r(n,q)=\frac1M
 \sum_{\substack{0\le a<M,\ (a,L_0)=1\\ \|a/M\|\ge 1/(nH)}}
 s_+(a/L_0)\,\frac{\log^r d(a/M)}{d(a/M)} .
\]
\begin{theorem}\label{th:main}
Uniformly for \(q\ge n^{40}\), as \(n=b^j\to\infty\),
\begin{align}
 S_0(n,q)&=-\frac{1-b^{-1}}{24\pi\Omega_+n^2}L(3)+O_{E,b,\Omega_+}(n^{-4}),\label{eq:zero}\\
 S_1(n,q)&=\frac{1-b^{-1}}{\pi\Omega_+}
 \left\{\frac{L''(1)}{2}+
 \frac{(\gamma+\log n)L(3)-L'(3)}{24n^2}\right\}
 +O_{E,b,\Omega_+}\bigl(n^{-4}(1+\log n)\bigr).\label{eq:one}
\end{align}
Here \(\gamma\) is Euler's constant. For each summand \(|1-\zeta_M^a|=d(a/M)\)
is a global algebraic unit.
\end{theorem}
The arguments \(\zeta_M^a\) here have order \(M\), since \(n=b^j\).
The identity \(\Phi_M(1)=1\) for \(M\) divisible by two distinct primes
gives the unit assertion, first for \(1-\zeta_M^a\), then for its absolute
value by squaring. Equation \eqref{eq:one} implies eventual nonvanishing
of \(S_1\) even when \(L''(1)=0\), since \(L(3)>0\) by the Euler product.
The distinction matters when using a nonzero finite sum to test a limiting value.

\section{Distribution and endpoints}
For \(\im z\ge0\) let \(J(z)=\sum a_n\e^{2\pi inz}/(2\pi i n^2)\).
The series converges absolutely and uniformly, by the bound
\(|a_n|\le d_n\sqrt n\) with \(d_n\) the divisor function.
Let \(G_\pm(x)\) be the corresponding real or imaginary part of
\(J(x)-J(0)\), divided by \(\Omega_\pm\). These continuous periodic
functions vanish at every integer. Write the distributions with these
primitives as \(\dd G_\pm\), and their dilates (formal densities
\(s_\pm(nx)\)) as \(\nu_{\pm,n}=\dd(G_\pm(nx)/n)\).
For smooth periodic tests the latter can equivalently be computed by its
Fourier series, using Abel limits. For a piecewise smooth test with a
cutoff we use integration by parts against the continuous primitive.

\begin{lemma}\label{lem:dist}
At fixed \(N\), on all intervals the normalized sum of the cusp symbols
\(s_\pm(a/h)\) differs from the increment of \(G_\pm\) by
\(O_{E,\Omega_\pm,\eps}(h^{-1/4+\eps})\). Near any integer,
\(G_\pm(x)=O_{E,\Omega_\pm}(\|x\|^2)\). In addition,
\(|s_\pm(a/h)|=O_{E,\Omega_\pm}(\log(2h))\).
\end{lemma}
\begin{proof}
The first assertion is Theorem 1.2 of Diamantis, Hoffstein, K{\i}ral and
Lee \cite{DHKL}, with fixed level and their imaginary symbol divided by
\(i\) for the minus sign. Endpoint conventions affect at most two atoms.
For the atom estimate, a continued-fraction path to a reduced cusp of
denominator at most \(h\) uses \(O(\log(2h))\) unimodular edges. The
integrals on these edges range over finitely many values since
\(\Gamma_0(N)\) has finite index. For the second assertion write
\(W_Nz=-1/(Nz)\). The Fricke identity gives
\(F(z)=F(0)+\lambda F(W_Nz)\) with \(\lambda=\pm1\).
In our normalization \(F(0)=L(1)=0\). An antiderivative of \(F(W_Nz)\)
based at zero is
\[
 Nz^2 J(W_Nz)-2N\int_0^z wJ(W_Nw)\,\dd w.
\]
Differentiation verifies this in the upper half-plane. It extends
continuously to the boundary and is \(O(|z|^2)\) there, since \(J\)
is bounded. Its product with \(\lambda\) is \(J(z)-J(0)\).
\end{proof}

Here and below the normalized sum for a level \(h\) means \(h^{-1}\sum\).
For a periodic piecewise smooth test \(K\), the discrete sum at \(nh\)
with symbols \(s_\pm(a/h)\) pairs the level-\(h\) samples with
\[
 {\cal R}_nK(y)=\frac1n\sum_{l=0}^{n-1}K((y+l)/n).
\]
The same folding holds for \(\nu_{\pm,n}\).
The circular variation satisfies \(\operatorname{Var}({\cal R}_nK)
\le\operatorname{Var}(K)/n\) by splitting \([0,1]\) into \(n\) pieces.
The sup norm is bounded by the variation plus the absolute mean of \(K\).
Consequently partial summation gives error
\begin{equation}\label{eq:fold}
 O_{E,\Omega_\pm,\eps}\bigl(h^{-1/4+\eps}
      [\operatorname{Var}(K)/n+|\textstyle\int_0^1 K|]\bigr).
\end{equation}
This estimate also holds for an explicitly specified inclusion of a cutoff
endpoint by the interval version of Lemma~\ref{lem:dist}.

\section{Two Mellin kernels}
Write \(\langle\nu,K\rangle\) for the distribution pairing.
For \(\re t\) large first define
\[
 I_n^+(t)=n^t\langle\nu_{+,n},d^{\,t-1}\rangle,\qquad
 I_n^-(t)=n^t\langle\nu_{-,n},\cot(\pi x)d(x)^t\rangle .
\]
Both pairings extend holomorphically to \(\re t>-1\) by integration
by parts at the endpoints, using Lemma~\ref{lem:dist} at fixed \(n\).
For example there the primitive is \(O(n x^2)\) and the differentiated
kernel is \(O_t(x^{\re t-2})\).

\begin{lemma}\label{lem:gamma}
On a fixed sufficiently small disk about zero, as identities of holomorphic
germs with locally uniform, differentiable remainders,
\begin{align}
 \Omega_+ I_n^+(t)&=C_+(t)\left[L(1+t)
 +\frac{t^3-t}{24 n^2}L(3+t)+O(t n^{-4})\right],\label{eq:plus}\\
 \Omega_- I_n^-(t)&=C_-(t)L(1+t)+O(n^{-2}),\label{eq:minus}\\
 C_+(t)&=\Gamma(t)\cos(\pi t/2)/\pi,\qquad
 C_-(t)=2\Gamma(1+t)\sin(\pi t/2)/(\pi t).\nonumber
\end{align}
\end{lemma}
\begin{proof}
The cosine coefficient of \(d^{\,t-1}\) at positive frequency \(l\),
with convention \(\int_0^1 K(x)\cos(2\pi lx)\,\dd x\), is
\[
 C_+(t)\frac{\Gamma(l+(1-t)/2)}{\Gamma(l+(1+t)/2)}.
\]
Indeed the beta integral at \(l=0\) gives the value there, and integration
by parts gives the ratio \((l+(1-t)/2)/(l+(1+t)/2)\) of consecutive
coefficients. The sine coefficient of \(\cot(\pi x)d^t\) is
\[
 C_-(t)\,l\frac{\Gamma(l-t/2)}{\Gamma(l+1+t/2)} .
\]
This follows by integrating \((d^t)'/(\pi t)\) against \(\sin(2\pi lx)\)
and applying the preceding cosine formula with exponent \(t\).
Multiply the coefficients at \(l=kn\) by \(n^t a_k/k\) and sum.
The ratio asymptotic in the even case, after removing \(l^{-t}\), is
\(1+(t^3-t)/(24l^2)+O(t l^{-4})\).
This follows by subtracting the two Stirling series for \(\log\Gamma\).
The complementary gamma arguments eliminate odd inverse powers, and
the residual term vanishes at \(t=0\); uniformity on a slightly larger
disk gives the factor \(t\). The odd ratio with its prefactor \(l\)
is \(l^{-t}(1+O_t(l^{-2}))\) by the same argument.
After subtraction the coefficient series converge normally near zero.
Starting from \(\re t\) large, continuation through the half-plane
\(\re t>-1/4\), with removable exceptional coefficient formulas read
by limits, proves \eqref{eq:plus} and \eqref{eq:minus}. Here finitely
many low frequencies may always be taken out of the Stirling expansion.
The pole of \(C_+\) at zero cancels under the vanishing hypotheses.
\end{proof}

To justify the passage to the finite squarefree sums uniformly, cut the
kernels off at \(\|x\|=1/(nH)\). For either of the shapes
\(d^{-1}\log^r d\) or \(\cot(\pi x)\log^r d\), \(0\le r\le2\),
the cut kernel has variation \(O_r(nH(1+\log(nH))^r)\) and absolute
integral \(O_r((1+\log(nH))^{r+1})\). On the removed intervals,
integration by parts and the bound \(O(n\|x\|^2)\) for the primitive
give \(O_r(H^{-1}(1+\log(nH))^r)\), including boundary terms.
Apply \eqref{eq:fold} at \(h=bq\) and at \(h=q\). Terms with \(b\mid a\) are
exactly \(1/b\) times the latter sum with the same dilation and cutoff.
The remaining terms with \(q\mid a,\ b\nmid a\) have symbols at level
\(b\); their normalized absolute contribution is
\(O_{E,b,r}((1+\log(n bq))^{r+1}/q)\) by the harmonic sum on the
\(nb\)-grid, with the origin omitted. These statements continue to
hold for general positive integers \(n\) if the sieve condition is
written \((a,bq)=1\).
In particular the discretization and sieve errors together, apart from
the multiplier \(1-b^{-1}\), are bounded by
\begin{equation}\label{eq:disc}
 O_{E,b,\Omega_\pm,r,\eps}\bigl(q^{-1/8+\eps}(1+\log(n bq))^{r+1}\bigr).
\end{equation}
For \(q\ge n^{40}\) this is \(O(n^{-4})\) with a sufficiently small
fixed \(\eps>0\), including every displayed logarithm. Finally expand
\(C_+(t)=(1/t-\gamma+O(t))/\pi\). The logarithm used here is \(\log d\),
so subtract \((\log n) I_n^+(0)\) from \(I_n^{+\,\prime}(0)\).
This proves Theorem~\ref{th:main}. Likewise, \eqref{eq:minus} and
\(C_-(0)=1\) give
\begin{equation}\label{eq:odd}
 \langle\nu_{-,n},\cot(\pi x)\log^2 d(x)\rangle
 =\frac{L''(1)}{\Omega_-}
 +O_{E,\Omega_-}\bigl(n^{-2}(1+\log n)^2\bigr).
\end{equation}

\section{A finite trace from rational points}
We give the precise algebraic hypotheses before passing to a family of
traces. Fix a real Weierstrass uniformization \(\C/\Lambda\simeq E(\C)\)
using a rational invariant differential. Take \(P_1,P_2,A\in E(\Q)^0\),
where the superscript denotes the real identity component, and assume
\(A,A+P_i\ne O\). Use real elliptic logarithms \(u_i,z_A\) for these
points \(P_i,A\). Choose a primitive purely imaginary period \(\omega\);
\(\eta_\omega=\zeta(z+\omega)-\zeta(z)\) is its quasi-period.
Here \(\zeta(z)\) is the Weierstrass zeta function. Put
\begin{equation}\label{eq:phase}
 Z_i=\exp\!\bigl(2\eta_\omega u_i
        -2\omega[\zeta(z_A+u_i)-\zeta(z_A)]\bigr).
\end{equation}
These are unit complex numbers. The product formulas below only use
their displayed values and make no independence assumption about them.

Let \(m\) be an odd prime and \(T=k\omega/m\pmod\Lambda\) with \(m\nmid k\).
Use a Weil function \(f_T\) with divisor \(m([T]-[O])\), and \(g_T\) with
divisor \([m]^*([T]-[O])\). Both exist over \(K_m=\Q(E[m])\).
Indeed their divisors are rational and principal over the algebraic closure;
the latter assertion follows from the elliptic group law, and descent
of the one-dimensional space gives rational functions over \(K_m\).
Then \(g_T^m/(f_T\circ[m])\) is constant. Choose real lifts \(B,R_i\)
under \([m]\) of \(A,P_i\) respectively. Put
\[
 h_i(T)=g_T(B+R_i)/g_T(B),\qquad v_i=h_i(T)/\overline{h_i(T)} .
\]
Choose admissible \(m,k\) so these evaluations have no zero or pole
(for nonzero real \(A,A+P_i\) and nonreal \(T\) this is automatic).
By choice of real lifts, complex conjugation replaces \(T\) by \(-T\)
in the evaluation ratio. Thus the \(v_i\) are algebraic and have
absolute value one. In a finite Galois extension containing the data,
over \(K_m\), translation of a lift by \(S\in E[m]\) multiplies \(g_T\)
by the Weil character \(e_m(S,T)\), up to a simultaneous inversion
convention. The translations of \(B\) cancel, leaving
\begin{equation}\label{eq:weil}
 \sigma(v_i)/v_i=e_m(\sigma R_i-R_i,T)^2,\qquad
 v_i^m=\frac{f_T(A+P_i)/f_T(A)}{f_{-T}(A+P_i)/f_{-T}(A)} .
\end{equation}
Changing the root \(v_i\) within \(\mu_m v_i\) preserves both identities.
For Weil functions and the translation convention see \cite{Weil}.

\begin{proposition}\label{prop:trace}
Fix \(L'\) divisible by \(L_0=bq\), put \(M=mL'\), and fix the embedding
in \(\C\). In an average over the finite Galois group fixing
\(K_m(\zeta_M)\), suppose the joint character in \eqref{eq:weil} has image
\(\mu_m^2\). For \(a\) with \(\zeta_{L'}^a v_i^m\ne1\),
\begin{equation}\label{eq:trace}
 \av_\sigma\prod_{i=1}^2\log|1-\zeta_M^a\sigma(v_i)|
 =m^{-2}\prod_{i=1}^2\log|1-\zeta_{L'}^a v_i^m|.
\end{equation}
Repeating one mark gives instead the average of the squares, which differs
from the square of its average by the variance over \(\mu_m\).
\end{proposition}
\begin{proof}
Uniformity on the joint image gives independence of the two coordinates.
For each coordinate take absolute values and logarithms in
\(\prod_{\xi\in\mu_m}(1-z\xi)=1-z^m\). With a repeated coordinate
the variance identity is the usual identity for the second moment.
\end{proof}
For non-CM \(E/\Q\) and independent \(P_1,P_2\) modulo torsion the
surjectivity hypothesis holds for every sufficiently large prime \(m\),
including after adjoining \(\zeta_M\).
Here are the details. For large \(m\) the residual image is
\(\mathrm{GL}_2(\mathbf F_m)\) by Serre \cite{Serre}, and the two
point classes are independent modulo \(m\) by finite generation.
Write \(V=E[m]\). Over \(K_m\) the Kummer image in \(V^2\) is a
\(\mathrm{GL}_2\)-submodule. A proper such submodule is killed by a
nonzero \(\mathbf F_m\)-linear combination \(V^2\to V\), since \(V^2\)
is semisimple and \(\operatorname{End}_{\mathrm{GL}_2}(V)=\mathbf F_m\).
The corresponding Kummer class would then be inflated from
\(H^1(\mathrm{GL}_2,V)=0\). This vanishing follows by applying the
cocycle identity to the central scalar \(2I\). This proves full image.
The intersection of the division extension of the marks with
\(K_m(\zeta_M)\) over \(K_m\) gives a quotient of \(V^2\) on which
the linear Galois group acts trivially, by abelianity of the cyclotomic
extension over \(\Q\). The scalar \(2I\) forces this quotient to vanish.
The Weil character for nonzero \(T\), squared at an odd prime, maps
\(V\) onto \(\mu_m\), proving the claim.

\section{Simultaneous convergence in a complex embedding and in the trace}
\begin{theorem}\label{th:phase}
Let \(E\) be non-CM and retain the order-of-vanishing hypothesis. For two
independent \(P_i\) as above there are integers \(k\to\infty\), primes
\(m\asymp k^3,\ q\asymp k^{1/2}\), and a power \(D\asymp k^{1/8}\)
of the fixed prime \(b\). Put \(L'=D bq,\ M=mL',\ H=\lceil q^{1/8}\rceil\).
There are choices of the roots in \eqref{eq:weil} for which
\begin{equation}\label{eq:tor}
 {\cal T}(v_1,v_2)=\frac1M
 \sum_{\substack{0\le a<M,\ (a,bq)=1\\ \|a/L'\|\ge 1/(DH)}}
 s_-(a/(bq))\cot(\pi a/M)
   \prod_{i=1}^2\log|1-\zeta_M^a v_i|
\end{equation}
and \(m^2\av_\sigma{\cal T}(\sigma v_1,\sigma v_2)\) both tend to
\((1-b^{-1})L''(1)/\Omega_-\).
Only \((a,bq)=1\) is imposed, so multiples of \(m\) are retained.
\end{theorem}
\begin{proof}
In the fixed uniformization a representative of \(f_T(z)\), up to a
scalar, is
\[
 \left(\frac{\sigma(z-k\omega/m)}{\sigma(z)}\right)^m
          \exp(k\eta_\omega z).
\]
It has the required divisor. Its periodicity follows from the sigma
quasi-periods and the Legendre relation. Using \(\sigma'/\sigma=\zeta\),
expansion of the logarithm of the ratios at \(z_A+u_i,z_A\) gives
\begin{equation}\label{eq:sigma}
 v_i^m=Z_i^k\exp(O_{E,A,P_i}(k^3/m^2)).
\end{equation}
Indeed subtracting the expansions at \(\pm k\omega/m\) leaves
the linear term in \eqref{eq:phase} and a cubic remainder. This describes the
third-kind differential behind the construction: \(\dd\log f_T\) has
simple poles of residues \(m,-m\) at \(T,O\); exponentiating its integral
from \(A\) to \(A+P_i\) gives the ratio in \eqref{eq:weil}.

Simultaneous Dirichlet approximation gives unbounded \(k\) with
\(\max_i|Z_i^k-1|\ll k^{-1/2}\). If the approximating integers stay
bounded along an unbounded approximation sequence, a common power is
exactly one and all its multiples work. Choose \(m\in[k^3,2k^3]\) and
\(q\in[k^{1/2},2k^{1/2}]\) prime by Bertrand's theorem, and choose the
largest power of \(b\) at most \(k^{1/8}\) for \(D\).
Finitely many forbidden primes are immaterial.
By \eqref{eq:sigma} the argument of each \(v_i^m\) can be chosen \(O(k^{-1/2})\).
Choose the corresponding \(m\)-th root in \eqref{eq:weil}. Then
\(|v_i-1|\ll k^{-1/2}/m\).

In the sum \eqref{eq:tor}, \(\|x\|\ge1/(mDH)\) for \(x=a/M\); on the corresponding
base grid, \(y=a/L'\) satisfies \(\|y\|\ge1/(DH)\). At either scale
the perturbation of the argument inside a logarithm is at most a constant
times \(k^{-1/2}DH\) relative to \(|1-\exp(2\pi ix)|\) or
\(|1-\exp(2\pi iy)|\), respectively. This quantity tends to zero
by a power of \(k\). The resulting change per logarithm is
\(O(k^{-1/2}DH)\); the other logarithm is \(O(\log(2M))\).
The normalized sum of absolute cotangents on the punctured grid is
\(O(\log(2M))\), and Lemma~\ref{lem:dist} bounds the symbols by
\(O(\log(2L_0))\). These estimates justify replacing the marks by \(1\)
in the chosen embedding and their powers by \(1\) on the base grid,
with errors \(o(1)\).

For the trace apply \eqref{eq:trace}. Each remaining factor, cutoff and modular symbol
is constant upon replacing \(a\) by \(a+hL'\), \(0\le h<m\).
The cotangent identity
\(\sum_{h=0}^{m-1}\cot(\pi(y+h)/m)=m\cot(\pi y)\)
therefore gives exactly the analogous base sum of modulus \(L'\),
with logarithm factors \(v_i^m\) and normalization \(1/L'\), after
multiplication of the trace by \(m^2\). The replacement just proved
turns the base sum into the plain odd quadrature at dilation \(D\).
Equations \eqref{eq:disc} and \eqref{eq:odd} give its limit.

For the chosen top embedding first consider the plain quadrature with
the single cutoff \(\|x\|\ge1/(mDH)\). This is \eqref{eq:disc}--\eqref{eq:odd} at dilation
\(mD\), and all logarithms of the scales are \(O(\log k)\). The actual
cutoff additionally removes intervals of radius \(1/(mDH)\) around
\(h/m,\ 1\le h<m\). In each such interval the cotangent is bounded
by \(O(m/\min(h,m-h))\). There are at most \(2L'/(DH)+2\)
grid points per interval and their logarithms are \(O(\log M)\).
Their total normalized contribution is
\(O((1/(DH)+1/L')\log^4(2M))=o(1)\).
This completes both limits.
\end{proof}

The same primes and torsion levels admit an additional comparison. Evaluate
\eqref{eq:weil} at the step \((k+1)\omega/m\) without changing other parameters.
The ratio of the two \(m\)-th powers in \eqref{eq:sigma} tends to \(Z_i\).
This is an embedding-dependent invariant of the displayed period
construction. The independence used in \eqref{eq:trace} comes instead from the Galois
image of the marked points. Neither finite-level independence nor equality
of the two limits determines whether \(L''(1)\) vanishes.

\subsection{A neighboring torsion direction}
For the algebraic Weierstrass invariants used here the unrotated phases
satisfy a stronger property. For every nonzero integer vector $(j_1,j_2)$,
the product $Z_1^{j_1}Z_2^{j_2}$ is transcendental. Indeed the zeta
addition theorem shows that
\[
 \zeta(z_A+u_i)-\zeta(z_A)-\zeta(u_i)\in\overline{\Q}.
\]
The addition term is a rational function of the algebraic Weierstrass
coordinates, with the duplication formula used when the two points
coincide. The hypotheses exclude a zero sum. Likewise, for
$u=j_1u_1+j_2u_2$ we have
$\zeta(u)-j_1\zeta(u_1)-j_2\zeta(u_2)\in\overline{\Q}$, by repeated
addition and oddness. In making the additions one first treats
multiples of each point and then adds the two resulting points;
independence eliminates zero intermediate sums. The point represented
by $u$ is algebraic and nontorsion. Consequently the product in
question is the inverse square of
\[
 \exp(\omega\zeta(u)-\eta_\omega u+\beta\omega)
\]
for an algebraic $\beta$. This number is transcendental by
Waldschmidt \cite[Thm.~7.2.1]{Wal}; the primitive imaginary period
satisfies $\omega/2\notin\Lambda$. In particular both phases differ
from one. Multiplicative independence here is not a statement about
algebraic independence of the phases.

One can observe the phase distinction in a second trace. In the
logarithms of \eqref{eq:tor} multiply every $v_i$ by the rational
scalar $\rho=1-M^{-1}$. For the aligned step the two limits already
proved are unchanged. Indeed $1-\rho^m\le1/L'$, so the added relative
perturbation at the cutoff on the lower mesh is $O(DH/L')=o(1)$
by a power of $k$, and the top mesh gives the same bound. Now replace
$k\omega/m$ by $(k+1)\omega/m$. On the same sequence the normalized
trace tends to zero.

To verify this, let $w_i=\rho^m v_i^m$ for the neighboring step.
Equation~\eqref{eq:sigma} gives $w_i\to Z_i$. Put
$B(y)=\prod_i\log|1-e^{2\pi iy}w_i|$ and $c=B(0)$.
The exact finite trace and cotangent folding reduce its normalized
value to the unit-filtered lower mesh with test $\cot(\pi y)B(y)$.
The term $c\cot(\pi y)$ tends to zero by the unlogged odd Mellin
formula at $t=0$ (its central term is proportional to $L(1)=0$,
and its finite-dilation error tends to zero); $c$ is bounded.
The function
$H_0(y)=(B(y)-B(0))\cot(\pi y)$ extends periodically across zero
and has bounded $C^1$ norm on a fixed arc there, as both $Z_i\ne1$.
On the complement its sup norm and variation are
$O_Z(\log^2 L')$. For each logarithm this follows by monotonicity
on the two semicircles relative to its peak and
$1-|w_i|\asymp1/L'$. Pairing $H_0$ with the distribution whose
primitive is $G_-(Dy)/D$ therefore costs $O_Z(\log^2 L'/D)$.
The mesh errors at levels $bq$ and $q$ are
$O_{E,Z,\epsilon}(q^{-1/4+\epsilon}\log^2L')$ by
\eqref{eq:fold}; multiples of $b$ are subtracted by exact folding,
and the remaining multiples of $q$ contribute
$O_{E,b,Z}(\log^2 L'/q)$ since their symbols have denominator at
most $b$. On the deleted endpoint interval $H_0$ is bounded and
the loss is $O((1/(DH)+1/L')\log(2L_0))$. Every error tends to zero.
The vanishing neighboring trace distinguishes torsion directions
when the aligned primitive derivative is nonzero. The limits may
both be zero.

\section{A real test with prescribed Kurihara residues}
For the final application we use $\Omega=\Omega_+$ and $m(x)=s_+(x)$, retaining the plus-symbol normalization above. The level and dilation variables in this section are independent of those in the previous construction.
Let $E/\Q$ have conductor $N$, and let $f(z)=\sum_{j\ge1}a_j e^{2\pi i jz}$ be its normalized newform. Fix $p\ge5$ and a positive real period $\Omega$ such that
\[
 m(x)=\frac{2\pi}{\Omega}\Re\left(i\int_{i\infty}^{x} f(z)\,dz\right)\in\Z_{(p)}
 \qquad(x\in\Q).
\]
This is a normalization condition on the plus modular-symbol lattice; the witness and the target use the same period. Assume $L(E,1)=0$ and root number $+1$. Then $L'(E,1)=0$. Choose distinct good primes $\ell_1,\ell_2\nmid pN$ with $\ell_i\equiv1$ and $a_{\ell_i}\equiv2\bmod p$. Put $n=\ell_1\ell_2$, $U_d=(\Z/d\Z)^\times$. Choose logarithms $\lambda_i:U_{\ell_i}\to\FF_p$ mapping the chosen generators to 1. Assume
\[
 w_0(b)=\lambda_1(b)\lambda_2(b),\qquad
 \delta=\sum_{b\in U_n}w_0(b)m(b/n)\ne0\quad\text{in }\FF_p.
\]
A nonzero value modulo a higher power of $p$ alone need not satisfy this condition. The discrete rank theorem has separate Selmer hypotheses \cite{Kim}. The companion note \cite{previous} treats the full-support finite-order central values. Our concern is the real limit on the base elliptic curve.

Take primes $q\nmid nNp$ satisfying
\begin{equation}\label{qcond}
 q\equiv1\bmod np,\qquad a_q\not\equiv2\bmod p .
\end{equation}
For the limiting statements we assume that these primes are infinite. A sufficient condition is that the residual image on $E[p]$ contains $\mathrm{SL}_2(\FF_p)$. Its commutator then contains this perfect subgroup, so $-I$ fixes the intersection of the division field with $\Q(\zeta_{np})$. Chebotarev applied to the compatible pair $(-I,1)$ gives positive density, since $-2\ne2\bmod p$.

\subsection{The Fricke kernel}
For a smooth periodic real even function $K$, put
\[
 \A_M(K)=\frac1M\sum_{a=0}^{M-1}m(a/M)K(a/M).
\]
Theorem 1.2 of \cite{DHKL}, followed by Stieltjes integration by parts, gives for $(M,N)=1$
\begin{align*}
 \A_M(K)&=\Lfun(K)+O_{f,\Omega,\epsilon}
       \big(M^{-1/4+\epsilon}(\|K\|_\infty+\|K'\|_1)\big),\\
 \Lfun(K)&=\frac1\Omega\sum_{j\ge1}\frac{a_j}{j}
       \int_0^1K(x)\cos(2\pi jx)\,dx .
\end{align*}
Indeed the limiting cumulative function is $(2\pi\Omega)^{-1}\sum_{j\ge1}a_j\sin(2\pi jx)/j^2$. This series converges absolutely, and the proof in \cite{DHKL} gives the displayed error uniformly in the endpoint (the translated cutoff phases have modulus one). The duplicated endpoint in their cumulative sum contributes zero since $m(0)=0$. In particular $\Lfun(\cos(2\pi j\cdot))=a_j/(2j\Omega)$.

An explicit real analytic test is
\[
 K_F(x)=8\pi\sum_{j\ge1}j\,\cos(2\pi jx)
       \int_{1/\sqrt N}^{\infty}e^{-2\pi jy}\log^2(\sqrt N y)\,dy .
\]
Every differentiated Fourier series converges uniformly on the real axis. The Mellin transform and the two lower vanishing derivatives give $L''(E,1)=2\pi\int_0^\infty f(iy)\log^2(\sqrt N y)\,dy$. The Fricke formula $f(i/(Ny))=Ny^2 f(iy)$ makes the latter integral twice its portion on $y\ge1/\sqrt N$. Hence $\Lfun(K_F)=L''(E,1)/\Omega$.

\subsection{Rounded quadratures}
For $M=nq$, set $w_M(a)=w_0(a\bmod n)$ on $U_M$ and zero elsewhere. In the residue class $w_M(a)\bmod p$, choose the nearest integer $k_a$ to $MK(a/M)$, breaking ties downward. Define
\[
 \mathcal R_M(K)=M^{-2}\sum_{a=0}^{M-1}k_a m(a/M).
\]
\begin{theorem}
For every fixed smooth real even $K$ and $q$ as in \eqref{qcond}, the rational number $\mathcal R_M(K)$ is a $p$-unit with residue $M^{-2}(a_q-2)\delta$. Along this set of primes it converges over $\mathbb R$ to $\Lfun(K)$. In particular $K_F$ gives the second central derivative divided by $\Omega$.
\end{theorem}
\begin{proof}
The one-prime Hecke norm relation \cite[\S3]{KN} gives for $b\in U_n$
\[
 \sum_{\substack{a\in U_{nq}\\a\equiv b\bmod n}}m(a/(nq))
       =a_qm(b/n)-m(qb/n)-m(q^{-1}b/n).
\]
The complete additive fiber gives $a_qm(b/n)-m(qb/n)$, and omitting the unique numerator divisible by $q$ subtracts $m(q^{-1}b/n)$, with inverse taken modulo $n$. Condition \eqref{qcond} reduces this to $(a_q-2)m(b/n)$. Reducing the integers $k_a$ gives the claimed residue. The continued fraction for a reduced fraction of denominator at most $M$ gives $O(\log(2M))$ unimodular segments. Their modular integrals take finitely many values since $\Gamma_0(N)$ has finite index in $\mathrm{SL}_2(\Z)$. Hence $|m(a/M)|\ll_{f,\Omega}\log(2M)$ and
\[
 |\mathcal R_M(K)-\A_M(K)|\ll_{f,\Omega}p\log(2M)/M .
\]
The fixed-denominator estimate finishes the proof.
\end{proof}

\paragraph{Bounded integer weights.}
One can avoid scaling the kernel by $M$. In increasing order of $0\le a<M$, let $k^-_a,k^+_a$ be the adjacent integers in the required residue class bracketing $K(a/M)$ (use the common value when it is already in the class). Choose one of them, $j_a$, so that $(j_a-K(a/M))m(a/M)$ has sign opposite the running error, allowing zero and choosing either sign if the previous error was zero. Write $D_b=\sum_{a<b}(j_a-K(a/M))m(a/M)$. Each possible step has magnitude at most $p|m(a/M)|$; by the choice of its sign, $|D_{a+1}|\le\max(|D_a|,p|m(a/M)|)$. Thus
\[
 |j_a|\le\|K\|_\infty+p,\qquad
 \max_{b\le M}|D_b|\ll_{E,\Omega}p\log(2M).
\]
The rational sum $M^{-1}\sum_{a<M}j_a m(a/M)$ consequently has real limit $\Lfun(K)$ with the same error bound and has nonzero residue $M^{-1}(a_q-2)\delta\bmod p$. Its reduced denominator is $O_{f,\Omega}(M)$. Indeed each normalized unimodular segment is a difference of two of the rational $p$-integral cusp symbols. There are only finitely many such segment values, so one fixed integer prime to $p$ clears all their denominators and those of every $m(a/M)$. The prefix bound is deterministic and does not require statistical independence of the symbols.

\subsection{Two zero-limit controls}
Choose a fixed integer lift $W$ of $w_0$ on $U_n$. The very same norm relation gives the exact identity
\[
 \frac1{nq}\sum_{a\in U_{nq}}W(a\bmod n)m(a/(nq))
       =\frac{a_q-2}{nq}\sum_{b\in U_n}W(b)m(b/n).
\]
Its left side is a rational $p$-unit at every admissible prime and is $O_{E,n,W}(q^{-1/2})$ by Hasse. In particular, it tends to zero without any rounding.

For an analytic control put $\ell=\ell_1$ and take
\[
 K_0(x)=a_\ell\cos(2\pi x)-\ell\cos(2\pi\ell x).
\]
Then $\Lfun(K_0)=0$. Nevertheless $\A_M(K_0)$ is nonzero and is a unit at some prime over $p$ of $\Q(\zeta_M)$. We prove this stronger assertion to distinguish the role of rounding from the arithmetic of the original modular-symbol sum.

Choose any order-$p$ character $\chi$ of $U_n$ nontrivial on both coordinates. In the group algebra of the quotient $C_p^2$, put $X_i=\sigma_i-1$ for the chosen generators. The projected theta element $\Theta_n=\sum_{b\in U_n}m(b/n)[b]$ satisfies
\[
 \Theta_n\equiv\delta X_1X_2\pmod{(p,(X_1,X_2)^3)}.
\]
The mixed coefficient is the defining moment. Specializing either variable to zero multiplies the smaller one-prime theta element by $a_{\ell_i}-\tau-\tau^{-1}$, where $\tau$ is the class of $\ell_i$ in the other factor. The multiplier has augmentation order at least two modulo $p$, while the one-prime element has zero augmentation by the same norm relation at level one and $m(0)=0$. This kills both pure terms. Since $X_i^p=0\bmod p$, no relations occur in degrees below three. Evaluating at $\bar\chi$ gives, with $v_p(p)=1$,
\[
 v_p(T_\chi)=2/(p-1),\qquad T_\chi=\sum_{b\in U_n}\bar\chi(b)m(b/n).
\]
Each $X_i$ evaluates to a unit multiple of $\zeta_p-1$, and $p$ has order $p-1>2$ in this uniformizer. See also \cite{previous}.

Write $\sigma_u(\zeta_M)=\zeta_M^u$ and $\tau(\chi)=\sum_{u\in U_n}\chi(u)e^{2\pi i u/n}$. The scalar character coefficients below are in the compositum with $\Q(\zeta_p)$. Then
\[
 \sum_{u\in U_M}\chi(u)\sigma_u\big(M\A_M(K_0)\big)
                =a_\ell\tau(\chi)(1+q-a_q)T_\chi .
\]
Indeed the $\ell$th frequency is killed by primitivity of $\chi$ at $\ell$. For the first frequency the $n$-Gauss factor is $\tau(\chi)\bar\chi(a)\chi(q)$ for $(a,n)=1$ and zero otherwise. The $q$-factor is $-1$ when $q\nmid a$ and $q-1$ when $q\mid a$. The cosine gives the same result as the positive exponential because $\chi$ is even. Equation \eqref{qcond} and the norm relation yield the displayed identity. Here $a_\ell$, $\tau(\chi)$ and $1+q-a_q$ are all units at $p$ (the squared complex norm of the Gauss sum is $n$, prime to $p$). Every conjugate of $M\A_M(K_0)$ lies in the $p$-integral part of a field unramified at $p$. If all had positive valuation they would all have valuation at least one, contrary to the right-hand side's valuation $2/(p-1)$. This proves the assertion at some place above $p$; it gives no lower bound at the distinguished real embedding.

Thus both the rational control and the analytic control converge to zero on the same curve. Inferring $L''(E,1)\ne0$ would require an estimate specific to $K_F$ beyond the residue calculation and the real quadrature error.


\begin{thebibliography}{9}\raggedright
\bibitem{DHKL} N. Diamantis, J. Hoffstein, E. M. K{\i}ral and M. Lee,
\emph{Additive twists and a conjecture by Mazur, Rubin and Stein},
J. Number Theory, Theorem 1.2.
\href{https://doi.org/10.1016/j.jnt.2019.11.016}{doi:10.1016/j.jnt.2019.11.016}.
\bibitem{Weil} A. Enge, \emph{Bilinear pairings on elliptic curves},
Enseign. Math. 61 (2015), 211--243, \S3.
\href{https://ems.press/content/serial-article-files/44297?nt=1}{Publisher's full text}.
\bibitem{Serre} J.-P. Serre, \emph{Propri\'et\'es galoisiennes des points
d'ordre fini des courbes elliptiques}, Invent. Math. 15 (1972), 259--331.
\href{https://doi.org/10.1007/BF01405086}{doi:10.1007/BF01405086}.
\bibitem{KN} D. Kriz and A. C. Nordentoft, \emph{Horizontal $p$-adic $L$-functions}, \S3. \href{https://arxiv.org/abs/2310.20678}{Preprint}.
\bibitem{Kim} C.-H. Kim, \emph{The refined Tamagawa number conjectures for $\mathrm{GL}_2$}, Corollary 1.5 and Conjecture 1.6. \href{https://arxiv.org/abs/2505.09121}{Preprint}.
\bibitem{previous} \emph{A Kurihara witness and the leading values of all full-support twists}, public project manuscript. \href{https://github.com/cosmosgeek0/BSD/blob/main/kurihara-full-support-twists.tex}{Source}.

\bibitem{Wal} M. Waldschmidt, \emph{Transcendence methods}, Queen's Papers in
Pure and Applied Mathematics 52 (1979), \S7.2, Theorem~7.2.1.
\href{https://webusers.imj-prg.fr/~michel.waldschmidt/articles/pdf/QueensPaper52.pdf}{Author's text}.
\end{thebibliography}
\end{document}

